\Chapter{30}

\Verse{1}
{
    \zA{话说林黛玉自与宝玉口角后也觉后悔，但又无去就他之理，因此日夜闷闷如有 所失。}
}
{
    \eA{Lin Tai-yü herself, for we will now resume our narrative, was also, ever
        since her tiff with Pao-yü, full of self-condemnation, yet as she did not
        see why she should run after him,}
}
{
}

\Verse{2}
{
    \zA{紫鹃也看出八九，便劝道：“论前儿的事，竟是姑娘太浮躁了些。}
}
{
    \eA{she continued, day and night, as despondent as she would have been had she
        lost some thing or other belonging to her. Tzu Chüan surmised her
        sentiments.}
}
{
}

\Verse{3}
{
    \zA{别人不知 宝玉的脾气，难道咱们也不知道？}
}
{
    \eA{"As regards what happened the other day," she advised her, "you were, after
        all, Miss, a little too hasty;}
}
{
}

\Verse{4}
{
    \zA{为那玉也不是闹了一遭两遭了。”}
}
{
    \eA{for if others don't understand that temperament of Pao-yü's, have you and I,
        surely, also no idea about it?}
}
{
}

\Verse{5}
{
    \zA{黛玉啐道：“呸! 你倒来替人派我的不是。}
}
{
    \eA{Besides, haven't there been already one or two rows on account of that very
        jade?" "Ts'ui!" exclaimed Tai-yü. "Have you come, on behalf of others,}
}
{
}

\Verse{6}
{
    \zA{我怎么浮躁了？”}
}
{
    \eA{to find fault with me? But how ever was I hasty?"}
}
{
}

\Verse{7}
{
    \zA{紫鹃笑道：“好好儿的，为什么铰了那 穗子？}
}
{
    \eA{"Why did you," smiled Tzu Chüan, "take the scissors and cut that tassel when
        there was no good reason for it?}
}
{
}

\Verse{8}
{
    \zA{不是宝玉只有三分不是，姑娘倒有七分不是？}
}
{
    \eA{So isn't Pao-yü less to blame than yourself, Miss? I've always found his
        behaviour towards you,}
}
{
}

\Verse{9}
{
    \zA{我看他素日在姑娘身上就好，皆 因姑娘小性儿，常要歪派他，才这么样。”}
}
{
    \eA{Miss, without a fault. It's all that touchy disposition of yours, which
        makes you so often perverse, that induces him to act as he does." Lin Tai-yü
        had every wish to make some suitable reply,}
}
{
}

\Verse{10}
{
    \zA{黛玉欲答话，只听院外叫门。}
}
{
    \eA{when she heard some one calling at the door. Tzu Chüan discerned the tone of
        voice.}
}
{
}

\Verse{11}
{
    \zA{紫鹃听了 听，笑道：“这是宝玉的声音，想必是来赔不是来了。”}
}
{
    \eA{"This sounds like Pao-yü's voice," she smiled. "I expect he's come to make
        his apologies." "I won't have any one open the door," Tai-yü cried at these
        words.}
}
{
}

\Verse{12}
{
    \zA{黛玉听了，说：“不许开 门！”}
}
{
    \eA{"Here you are in the wrong again, Miss," Tzu Chüan observed.}
}
{
}

\Verse{13}
{
    \zA{紫鹃道：“姑娘又不是了，这么热天，毒日头地下，晒坏了他，如何使得呢。”}
}
{
    \eA{"How will it ever do to let him get a sunstroke and come to some harm on a
        day like this, and under such a scorching sun?" Saying this, she speedily
        walked out and opened the door.}
}
{
}

\Verse{14}
{
    \zA{口里说着，便出去开门，果然是宝玉。}
}
{
    \eA{It was indeed Pao-yü. While ushering him in, she gave him a smile. "I
        imagined,"}
}
{
}

\Verse{15}
{
    \zA{一面让他进来，一面笑着说道：“我只当宝 二爷再不上我们的门了，谁知道这会子又来了。”}
}
{
    \eA{she said, "that you would never again put your foot inside our door, Master
        Secundus. But here you are once more and quite unexpectedly!" "You have by
        dint of talking," Pao-yü laughed, "made much ado of nothing;}
}
{
}

\Verse{16}
{
    \zA{宝玉笑道：“你们把极小的事倒 说大了，好好的为什么不来？}
}
{
    \eA{and why shouldn't I come, when there's no reason for me to keep away? Were I
        even to die, my spirit too will come a hundred times a day! But is cousin
        quite well?"}
}
{
}

\Verse{17}
{
    \zA{我就死了，魂也要一日来一百遭。}
}
{
    \eA{"She is," replied Tzu Chüan, "physically all right; but, mentally, her
        resentment is not quite over."}
}
{
}

\Verse{18}
{
    \zA{妹妹可大好了？”}
}
{
    \eA{"I understand," continued Pao-yü with a smile.}
}
{
}

\Verse{19}
{
    \zA{紫鹃道：“身上病好了，只是心里气还不大好。”}
}
{
    \eA{"But resentment, for what?" With this inquiry, he wended his steps inside
        the apartment. He then caught sight of Lin Tai-yü reclining on the bed in
        the act of crying.}
}
{
}

\Verse{20}
{
    \zA{宝玉笑道：“我知道了，有什么 气呢。”}
}
{
    \eA{Tai-yü had not in fact shed a tear, but hearing Pao-yü break in upon her,
        she could not help feeling upset.}
}
{
}

\Verse{21}
{
    \zA{一面说着，一面进来。}
}
{
    \eA{She found it impossible therefore to prevent her tears from rolling down her
        cheeks.}
}
{
}

\Verse{22}
{
    \zA{只见黛玉又在床上哭。}
}
{
    \eA{Pao-yü assumed a smiling expression and drew near the bed.}
}
{
}

\Verse{23}
{
    \zA{那黛玉本不曾哭，听见宝玉来，由不得伤心，止不住滚下泪来。}
}
{
    \eA{"Cousin, are you quite well again?" he inquired. Tai-yü simply went on
        drying her tears, and made no reply of any kind. Pao-yü approached the bed,
        and sat on the edge of it.}
}
{
}

\Verse{24}
{
    \zA{宝玉笑着走近 床来道：“妹妹身上可大好了？”}
}
{
    \eA{"I know," he smiled, "that you're not vexed with me. But had I not come,
        third parties would have been allowed to notice my absence,}
}
{
}

\Verse{25}
{
    \zA{黛玉只顾拭泪，并不答应。}
}
{
    \eA{and it would have appeared to them as if we had had another quarrel.}
}
{
}

\Verse{26}
{
    \zA{宝玉因便挨在床沿上 坐了，一面笑道：“我知道你不恼我，但只是我不来，叫旁人看见，倒像是咱们又 拌了嘴的似的。}
}
{
    \eA{And had I to wait until they came to reconcile us, would we not by that time
        become perfect strangers? It would be better, supposing you wish to beat me
        or blow me up, that you should please yourself and do so now; but whatever
        you do, don't give me the cold shoulder!"}
}
{
}

\Verse{27}
{
    \zA{要等他们来劝咱们，那时候儿岂不咱们倒觉生分了？}
}
{
    \eA{Continuing, he proceeded to call her "my dear cousin" for several tens of
        times. Tai-yü had resolved not to pay any more heed to Pao-yü. When she,
        however, now heard Pao-yü urge:}
}
{
}

\Verse{28}
{
    \zA{不如这会子你 要打要骂，凭你怎么样，千万别不理我！”}
}
{
    \eA{"don't let us allow others to know anything about our having had a quarrel,
        as it will look as if we had become thorough strangers," it once more became
        evident to her,}
}
{
}

\Verse{29}
{
    \zA{说着，又把“好妹妹”叫了几十声。}
}
{
    \eA{from this single remark, that she was really dearer and nearer to him than
        any of the other girls,}
}
{
}

\Verse{30}
{
    \zA{黛 玉心里原是再不理宝玉的，这会子听见宝玉说“别叫人知道咱们拌了嘴就生分了似 的”这一句话，又可见得比别人原亲近，因又掌不住，便哭道：“你也不用来哄我!
        从今以后，我也不敢亲近二爷，权当我去了。”}
}
{
    \eA{so she could not refrain from saying sobbingly: "You needn't have come to
        chaff me! I couldn't presume henceforward to be on friendly terms with you,
        Master Secundus! You should treat me as if I were gone!" At these words,
        Pao-yü gave way to laughter. "Where are you off to?" he inquired. "I'm going
        back home," answered Tai-yü.}
}
{
}

\Verse{31}
{
    \zA{宝玉听了笑道：“你往那里去呢？”}
}
{
    \eA{"I'll go along with you then," smiled Pao-yü. "But if I die?" asked Tai-yü.
        "Well, if you die,"}
}
{
}

\Verse{32}
{
    \zA{黛玉道：“我回家去。”}
}
{
    \eA{rejoined Pao-yü, "I'll become a bonze." The moment Tai-yü caught this reply,}
}
{
}

\Verse{33}
{
    \zA{宝玉笑道：“我跟了去。”}
}
{
    \eA{she hung down her head. "You must, I presume, be bent upon dying?"}
}
{
}

\Verse{34}
{
    \zA{黛玉道：“我死了呢？”}
}
{
    \eA{she cried. "But what stuff and nonsense is this you're talking?}
}
{
}

\Verse{35}
{
    \zA{宝玉 道：“你死了，我做和尚。”}
}
{
    \eA{You've got so many beloved elder and younger cousins in your family, and how
        many bodies will you have to go and become bonzes,}
}
{
}

\Verse{36}
{
    \zA{黛玉一闻此言，登时把脸放下来，问道：“想是你要 死了!胡说的是什么？}
}
{
    \eA{when by and bye they all pass away! But to-morrow I'll tell them about this
        to judge for themselves what your motives are!" Pao-yü was himself aware of
        the fact that this rejoinder had been recklessly spoken,}
}
{
}

\Verse{37}
{
    \zA{你们家倒有几个亲姐姐亲妹妹呢!明儿都死了，你几个身子做 和尚去呢？}
}
{
    \eA{and he was seized with regret. His face immediately became suffused with
        blushes. He lowered his head and had not the courage to utter one word more.
        Fortunately, however, there was no one present in the room.}
}
{
}

\Verse{38}
{
    \zA{等我把这个话告诉别人评评理。”}
}
{
    \eA{Tai-yü stared at him for ever so long with eyes fixed straight on him, but
        losing control over her temper,}
}
{
}

\Verse{39}
{
    \zA{宝玉自知说的造次了，后悔不来，登 时脸上红涨，低了头不敢作声。}
}
{
    \eA{"Ai!" she shouted, "can't you speak?" Then when she perceived Pao-yü reduced
        to such straits as to turn purple, she clenched her teeth and spitefully
        gave him, on the forehead, a fillip with her finger.}
}
{
}

\Verse{40}
{
    \zA{幸而屋里没人。}
}
{
    \eA{"Heug!" she cried gnashing her teeth,}
}
{
}

\Verse{41}
{
    \zA{黛玉两眼直瞪瞪的瞅了他半天，气的“嗳”了一声，说不出话来。}
}
{
    \eA{"you, this……" But just as she had pronounced these two words, she heaved
        another sigh, and picking up her handkerchief, she wiped her tears. Pao-yü
        treasured at one time numberless tender things in his mind, which he meant
        to tell her,}
}
{
}

\Verse{42}
{
    \zA{见宝玉别的 脸上紫涨，便咬着牙，用指头狠命的在他额上戳了一下子，“哼”了一声，说道： “你这个――”刚说了三个字，便又叹了一口气，仍拿起绢子来擦眼泪。}
}
{
    \eA{but feeling also, while he smarted under the sting of self-reproach (for the
        indiscretion he had committed), Tai-yü give him a rap, he was utterly
        powerless to open his lips, much though he may have liked to speak, so he
        kept on sighing and snivelling to himself. With all these things therefore
        to work upon his feelings, he unwillingly melted into tears. He tried to
        find his handkerchief to dry his face with,}
}
{
}

\Verse{43}
{
    \zA{宝玉心里 原有无限的心事，又兼说错了话，正自后悔；又见黛玉戳他一下子，要说也说不出 来，自叹自泣：因此自己也有所感，不觉掉下泪来。}
}
{
    \eA{but unexpectedly discovering that he had again forgotten to bring one with
        him, he was about to make his coat-sleeve answer the purpose, when Tai-yü,
        albeit her eyes were watery, noticed at a glance that he was going to use
        the brand-new coat of grey coloured gauze he wore, and while wiping her own,
        she turned herself round, and seized a silk kerchief thrown over the pillow,}
}
{
}

\Verse{44}
{
    \zA{要用绢子揩拭，不想又忘了带 来，便用衫袖去擦。}
}
{
    \eA{and thrust it into Pao-yü's lap. But without saying a word, she screened her
        face and continued sobbing. Pao-yü saw the handkerchief she threw, and
        hastily snatching it,}
}
{
}

\Verse{45}
{
    \zA{黛玉虽然哭着，却一眼看见他穿着簇新藕合纱衫，竟去拭泪， 便一面自己拭泪，一面回身将枕上搭的一方绡帕拿起来向宝玉怀里一摔，一语不 发，仍掩面而泣。}
}
{
    \eA{he wiped his tears. Then drawing nearer to her, he put out his hand and
        clasped her hand in his, and smilingly said to her: "You've completely
        lacerated my heart, and do you still cry? But let's go; I'll come along with
        you and see our venerable grandmother." Tai-yü thrust his hand aside.}
}
{
}

\Verse{46}
{
    \zA{宝玉见他摔了帕子来，忙接住拭了泪，又挨近前些，伸手拉了他 一只手，笑道：“我的五脏都揉碎了，你还只是哭。}
}
{
    \eA{"Who wants to go hand in hand with you?" she cried. "Here we grow older day
        after day, but we're still so full of brazen-faced effrontery that we don't
        even know what right means?" But scarcely had she concluded before she heard
        a voice say aloud: "They're all right!" Pao-yü and Tai-yü were little
        prepared for this surprise,}
}
{
}

\Verse{47}
{
    \zA{走罢，我和你到老太太那里去 罢。”}
}
{
    \eA{and they were startled out of their senses. Turning round to see who it was,
        they caught sight of lady Feng running in,}
}
{
}

\Verse{48}
{
    \zA{黛玉将手一摔道：“谁和你拉拉扯扯的!一天大似一天，还这么涎皮赖脸的， 连个理也不知道。”}
}
{
    \eA{laughing and shouting. "Our old lady," she said, "is over there, giving way
        to anger against heaven and earth. She would insist upon my coming to find
        out whether you were reconciled or not. 'There's no need for me to go and
        see,'}
}
{
}

\Verse{49}
{
    \zA{一句话没说完，只听嚷道：“好了！”}
}
{
    \eA{I told her, 'they will before the expiry of three days, be friends again of
        their own accord.'}
}
{
}

\Verse{50}
{
    \zA{宝黛两个不防，都唬了一跳。}
}
{
    \eA{Our venerable ancestor, however, called me to account, and maintained that I
        was lazy;}
}
{
}

\Verse{51}
{
    \zA{回头看时， 只见凤姐儿跑进来，笑道：“老太太在那里抱怨天，抱怨地，只叫我来瞧瞧你们好 了没有，我说：‘不用瞧，过不了三天，他们自己就好了。’}
}
{
    \eA{so here I come! But my words have in very deed turned out true. I don't see
        why you two should always be wrangling! For three days you're on good terms
        and for two on bad. You become more and more like children. And here you are
        now hand in hand blubbering! But why did you again yesterday become like
        black-eyed fighting cocks? Don't you yet come with me to see your
        grandmother and make an old lady like her set her mind at ease a bit?"}
}
{
}

\Verse{52}
{
    \zA{老太太骂我，说我懒； 我来了，果然应了我的话了。}
}
{
    \eA{While reproaching them, she clutched Tai-yü's hand and was trudging away,
        when Tai-yü turned her head round and called out for her servant-girls.}
}
{
}

\Verse{53}
{
    \zA{――也没见你们两个!有些什么可拌的，三日好了， 两日恼了，越大越成了孩子了。}
}
{
    \eA{But not one of them was in attendance. "What do you want them for again?"
        lady Feng asked. "I am here to wait on you!" Still speaking, she pulled her
        along on their way, with Pao-yü following in their footsteps.}
}
{
}

\Verse{54}
{
    \zA{有这会子拉着手哭的，昨儿为什么又成了‘乌眼鸡’ 似的呢？}
}
{
    \eA{Then making their exit out of the garden gate, they entered dowager lady
        Chia's suite of rooms. "I said that it was superfluous for any one to
        trouble,"}
}
{
}

\Verse{55}
{
    \zA{还不跟着我到老太太跟前，叫老人家也放点儿心呢。”}
}
{
    \eA{lady Feng smiled, "as they were sure of themselves to become reconciled; but
        you, dear ancestor, so little believed it that you insisted upon my going to
        act the part of mediator.}
}
{
}

\Verse{56}
{
    \zA{说着，拉了黛玉就 走。}
}
{
    \eA{Yet when I got there, with the intention of inducing them to make it up,}
}
{
}

\Verse{57}
{
    \zA{黛玉回头叫丫头们，一个也没有。}
}
{
    \eA{I found them, though one did not expect it, in each other's company,
        confessing their faults,}
}
{
}

\Verse{58}
{
    \zA{凤姐道：“又叫他们做什么，有我伏侍呢。”}
}
{
    \eA{and laughing and chatting. Just like a yellow eagle clutching the feet of a
        kite were those two hanging on to each other.}
}
{
}

\Verse{59}
{
    \zA{一面说，一面拉着就走，宝玉在后头跟着。}
}
{
    \eA{So where was the necessity for any one to go?" These words evoked laughter
        from every one in the room.}
}
{
}

\Verse{60}
{
    \zA{出了园门，到了贾母跟前，凤姐笑道： “我说他们不用人费心，自己就会好的，老祖宗不信，一定叫我去说和。}
}
{
    \eA{Pao-ch'ai, however, was present at the time so Lin Tai-yü did not retort,
        but went and ensconced herself in a seat near her grandmother. When Pao-yü
        noticed that no one had anything to say, he smilingly addressed himself to
        Pao-ch'ai. "On cousin Hsüeh P'an's birth-day,"}
}
{
}

\Verse{61}
{
    \zA{赶我到那 里说和，谁知两个人在一块儿对赔不是呢，倒像‘黄鹰抓住鹞子的脚’――两个人 都‘扣了环’了!那里还要人去说呢？”}
}
{
    \eA{he remarked, "I happened again to be unwell, so not only did I not send him
        any presents, but I failed to go and knock my head before him. Yet cousin
        knows nothing about my having been ill, and it will seem to him that I had
        no wish to go, and that I brought forward excuses so as to avoid paying him
        a visit.}
}
{
}

\Verse{62}
{
    \zA{说的满屋里都笑起来。}
}
{
    \eA{If to-morrow you find any leisure, cousin, do therefore explain matters for
        me to him."}
}
{
}

\Verse{63}
{
    \zA{此时宝钗正在这里，那黛玉只一言不发，挨着贾母坐下。}
}
{
    \eA{"This is too much punctiliousness!" smiled Pao-ch'ai. "Even had you insisted
        upon going, we wouldn't have been so arrogant as to let you put yourself to
        the trouble,}
}
{
}

\Verse{64}
{
    \zA{宝玉没什么说的，便 向宝钗笑道：“大哥哥好日子，偏我又不好，没有别的礼送，连个头也不磕去。}
}
{
    \eA{and how much less when you were not feeling well? You two are cousins and
        are always to be found together the whole day; if you encourage such ideas,
        some estrangement will, after all, arise between you." "Cousin," continued
        Pao-yü smilingly,}
}
{
}

\Verse{65}
{
    \zA{大 哥哥不知道我病，倒像我推故不去似的。}
}
{
    \eA{"you know what to say; and so long as you're lenient with me all will be all
        right. But how is it," he went on to ask, "that you haven't gone over to see
        the theatricals?"}
}
{
}

\Verse{66}
{
    \zA{倘或明儿姐姐闲了，替我分辩分辩。”}
}
{
    \eA{"I couldn't stand the heat" rejoined Pao-ch'ai. "I looked on while two plays
        were being sung,}
}
{
}

\Verse{67}
{
    \zA{宝 钗笑道：“这也多事。}
}
{
    \eA{but I found it so intensely hot, that I felt anxious to retire.}
}
{
}

\Verse{68}
{
    \zA{你就要去，也不敢惊动，何况身上不好。}
}
{
    \eA{But the visitors not having dispersed, I had to give as an excuse that I
        wasn't feeling up to the mark,}
}
{
}

\Verse{69}
{
    \zA{弟兄们常在一处， 要存这个心倒生分了。”}
}
{
    \eA{and so came away at once." Pao-yü, at these words, could not but feel ill at
        ease. All he could do was to feign another smile.}
}
{
}

\Verse{70}
{
    \zA{宝玉又笑道：“姐姐知道体谅我就好了。”}
}
{
    \eA{"It's no wonder," he observed, "that they compare you, cousin, to Yang
        Kuei-fei; for she too was fat and afraid of hot weather."}
}
{
}

\Verse{71}
{
    \zA{又道：“姐姐 怎么不听戏去？”}
}
{
    \eA{Hearing this, Pao-ch'ai involuntarily flew into a violent rage. Yet when
        about to call him to task,}
}
{
}

\Verse{72}
{
    \zA{宝钗道：“我怕热。}
}
{
    \eA{she found that it would not be nice for her to do so.}
}
{
}

\Verse{73}
{
    \zA{听了两出，热的很，要走呢，客又不散；我 少不得推身上不好，就躲了。”}
}
{
    \eA{After some reflection, the colour rushed to her cheeks. Smiling ironically
        twice, "I may resemble," she said, "Yang Kuei-fei, but there's not one of
        you young men, whether senior or junior,}
}
{
}

\Verse{74}
{
    \zA{宝玉听说，自己由不得脸上没意思，只得又搭讪笑 道：“怪不得他们拿姐姐比杨妃，原也富胎些。”}
}
{
    \eA{good enough to play the part of Yang Kuo-chung." While they were bandying
        words, a servant-girl Ch'ing Erh, lost sight of her fan and laughingly
        remarked to Pao-ch'ai: "It must be you, Miss Pao, who have put my fan away
        somewhere or other;}
}
{
}

\Verse{75}
{
    \zA{宝钗听说，登时红了脸，待要发 作，又不好怎么样；回思了一回，脸上越下不来，便冷笑了两声，说道：“我倒像 杨妃，只是没个好哥哥好兄弟可以做得杨国忠的！”}
}
{
    \eA{dear mistress, do let me have it!" "You'd better be mindful!" rejoined
        Pao-ch'ai, shaking her finger at her. "With whom have I ever been up to
        jokes, that you come and suspect me? Have I hitherto laughed and smirked
        with you?}
}
{
}

\Verse{76}
{
    \zA{正说着，可巧小丫头靓儿因不 见了扇子，和宝钗笑道：“必是宝姑娘藏了我的。}
}
{
    \eA{There's that whole lot of girls, go and ask them about it!" At this
        suggestion, Ch'ing Erh made her escape. The consciousness then burst upon
        Pao-yü, that he had again been inconsiderate in his speech,}
}
{
}

\Verse{77}
{
    \zA{好姑娘，赏我罢。”}
}
{
    \eA{in the presence of so many persons,}
}
{
}

\Verse{78}
{
    \zA{宝钗指着他 厉声说道：“你要仔细!你见我和谁玩过!有和你素日嘻皮笑脸的那些姑娘们，你该 问他们去！”}
}
{
    \eA{and he was overcome by a greater sense of shame than when, a short while
        back, he had been speaking with Lin Tai-yü. Precipitately turning himself
        round, he went, therefore, and talked to the others as well. The sight of
        Pao-yü poking fun at Pao-ch'ai gratified Tai-yü immensely.}
}
{
}

\Verse{79}
{
    \zA{说的靓儿跑了。}
}
{
    \eA{She was just about to put in her word and also seize the opportunity of
        chaffing her,}
}
{
}

\Verse{80}
{
    \zA{宝玉自知又把话说造次了，当着许多人，比才在黛玉 跟前更不好意思，便急回身，又向别人搭讪去了。}
}
{
    \eA{but as Ch'ing Erh unawares asked for her fan and Pao-ch'ai added a few more
        remarks, she at once changed her purpose. "Cousin Pao-ch'ai," she inquired,
        "what two plays did you hear?" Pao-ch'ai caught the expression of
        gratification in Tai-yü's countenance,}
}
{
}

\Verse{81}
{
    \zA{黛玉听见宝玉奚落宝钗，心中着实得意，才要搭言，也趁势取个笑儿，不想靓 儿因找扇子，宝钗又发了两句话，他便改口说道：“宝姐姐，你听了两出什么戏？”}
}
{
    \eA{and concluded that she had for a certainty heard the raillery recently
        indulged in by Pao-yü and that it had fallen in with her own wishes; and
        hearing her also suddenly ask the question she did, she answered with a
        significant laugh: "What I saw was: 'Li Kuei blows up Sung Chiang and
        subsequently again tenders his apologies'." Pao-yü smiled. "How is it,"}
}
{
}

\Verse{82}
{
    \zA{宝钗因见黛玉面上有得意之态，一定是听了宝玉方才奚落之言，遂了他的心愿。}
}
{
    \eA{he said, "that with such wide knowledge of things new as well as old; and
        such general information as you possess, you aren't even up to the name of a
        play, and that you've come out with such a whole string of words.}
}
{
}

\Verse{83}
{
    \zA{忽 又见他问这话，便笑道：“我看的是李逵骂了宋江，后来又赔不是。”}
}
{
    \eA{Why, the real name of the play is: 'Carrying a birch and begging for
        punishment'". "Is it truly called: 'Carrying a birch and begging for
        punishment'"? Pao-ch'ai asked with laugh.}
}
{
}

\Verse{84}
{
    \zA{宝玉便笑道： “姐姐通今博古，色色都知道，怎么连这一出戏的名儿也不知道，就说了这么一套。}
}
{
    \eA{"But you people know all things new and old so are able to understand the
        import of 'carrying a birch and begging for punishment.' As for me I've no
        idea whatever what 'carrying a birch and begging for punishment' implies."
        One sentence was scarcely ended when Pao-yü and Tai-yü felt guilty in their
        consciences; and by the time they heard all she said,}
}
{
}

\Verse{85}
{
    \zA{这叫做《负荆请罪》。”}
}
{
    \eA{they were quite flushed from shame. Lady Feng did not, it is true,}
}
{
}

\Verse{86}
{
    \zA{宝钗笑道：“原来这叫‘负荆请罪’!你们通今博古，才 知道‘负荆请罪’，我不知什么叫‘负荆请罪’。”}
}
{
    \eA{fathom the gist of what had been said, but at the sight of the expression
        betrayed on the faces of the three cousins, she readily got an inkling of
        it. "On this broiling hot day," she inquired laughing also; "who still eats
        raw ginger?"}
}
{
}

\Verse{87}
{
    \zA{一句话未说了，宝玉黛玉二人 心里有病，听了这话，早把脸羞红了。}
}
{
    \eA{None of the party could make out the import of her insinuation. "There's no
        one eating raw ginger," they said. Lady Feng intentionally then brought her
        hands to her cheeks,}
}
{
}

\Verse{88}
{
    \zA{凤姐这些上虽不通，但只看他三人的形景， 便知其意，也笑问道：“这们大热的天，谁还吃生姜呢？”}
}
{
    \eA{and rubbing them, she remarked with an air of utter astonishment, "Since
        there's no one eating raw ginger, how is it that you are all so fiery in the
        face?" Hearing this, Pao-yü and Tai-yü waxed more uncomfortable than ever.
        So much so, that Pao-ch'ai,}
}
{
}

\Verse{89}
{
    \zA{众人不解，便道：“没 有吃生姜的。”}
}
{
    \eA{who meant to continue the conversation, did not think it nice to say
        anything more when she saw how utterly abashed Pao-yü was and how changed
        his manner.}
}
{
}

\Verse{90}
{
    \zA{凤姐故意用手摸着腮，诧异道：“既没人吃生姜，怎么这么辣辣的 呢？”}
}
{
    \eA{Her only course was therefore to smile and hold her peace. And as the rest
        of the inmates had not the faintest notion of the drift of the remarks
        exchanged between the four of them,}
}
{
}

\Verse{91}
{
    \zA{宝玉黛玉二人听见这话，越发不好意思了。}
}
{
    \eA{they consequently followed her lead and put on a smile. In a short while,
        however, Pao-ch'ai and lady Feng took their leave.}
}
{
}

\Verse{92}
{
    \zA{宝钗再欲说话，见宝玉十分羞愧， 形景改变，也就不好再说，只得一笑收住。}
}
{
    \eA{"You've also tried your strength with them," Tai-yü said to Pao-yü
        laughingly. "But they're far worse than I. Is every one as simple in mind
        and dull of tongue as I am as to allow people to say whatever they like."}
}
{
}

\Verse{93}
{
    \zA{别人总没解过他们四个人的话来，因此 付之一笑。}
}
{
    \eA{Pao-yü was inwardly giving way to that unhappiness, which had been
        occasioned by Pao-ch'ai's touchiness, so when he also saw Tai-yü approach
        him and taunt him,}
}
{
}

\Verse{94}
{
    \zA{一时宝钗凤姐去了，黛玉向宝玉道：“你也试着比我利害的人了。}
}
{
    \eA{displeasure keener than ever was aroused in him. A desire then asserted
        itself to speak out his mind to her, but dreading lest Tai-yü should he in
        one of her sensitive moods,}
}
{
}

\Verse{95}
{
    \zA{谁都像我心 拙口夯的，由着人说呢！”}
}
{
    \eA{he, needless to say, stifled his anger and straightway left the apartment in
        a state of mental depression. It happened to be the season of the greatest
        heat.}
}
{
}

\Verse{96}
{
    \zA{宝玉正因宝钗多心，自己没趣儿，又见黛玉问着他，越 发没好气起来。}
}
{
    \eA{Breakfast time too was already past, and masters as well as servants were,
        for the most part, under the influence of the lassitude felt on lengthy
        days. As Pao-yü therefore strolled,}
}
{
}

\Verse{97}
{
    \zA{欲待要说两句，又怕黛玉多心，说不得忍气，无精打彩，一直出来。}
}
{
    \eA{from place to place, his hands behind his back he heard not so much as the
        caw of a crow. Issuing out of his grandmother's compound on the near side,
        he wended his steps westwards,}
}
{
}

\Verse{98}
{
    \zA{谁知目今盛暑之际，又当早饭已过，各处主仆人等多半都因日长神倦，宝玉背 着手，到一处，一处鸦雀无声。}
}
{
    \eA{and crossed the passage, on which lady Feng's quarters gave. As soon as he
        reached the entrance of her court, he perceived the door ajar. But aware of
        lady Feng's habit of taking, during the hot weather, a couple of hours'
        siesta at noon,}
}
{
}

\Verse{99}
{
    \zA{从贾母这里出来往西，走过了穿堂便是凤姐的院落。}
}
{
    \eA{he did not feel it a convenient moment to intrude. Walking accordingly
        through the corner door, he stepped into Madame Wang's apartment.}
}
{
}

\Verse{100}
{
    \zA{到他院门前，只见院门掩着，知道凤姐素日的规矩，每到天热，午间要歇一个时辰 的，进去不便。}
}
{
    \eA{Here he discovered several waiting-maids, dosing with their needlework
        clasped in their hands. Madame Wang was asleep on the cool couch in the
        inner rooms. Chin Ch'uan-erh was sitting next to her massaging her legs.}
}
{
}

\Verse{101}
{
    \zA{遂进角门，来到王夫人上房里。}
}
{
    \eA{But she too was quite drowsy, and her eyes wore all awry. Pao-yü drew up to
        her with gentle tread.}
}
{
}

\Verse{102}
{
    \zA{只见几个丫头手里拿着针线，却打 盹儿。}
}
{
    \eA{The moment, however, that he unfastened the pendants from the earrings she
        wore, Chin Ch'uan opened her eyes,}
}
{
}

\Verse{103}
{
    \zA{王夫人在里间凉床上睡着，金钏儿坐在傍边捶腿，也乜斜着眼乱恍。}
}
{
    \eA{and realised that it was no one than Pao-yü. "Are you feeling so worn out!"
        he smilingly remarked in a low tone of voice. Chin Ch'uan pursed up her lips
        and gave him a smile.}
}
{
}

\Verse{104}
{
    \zA{宝玉轻 轻的走到跟前，把他耳朵上的坠子一摘。}
}
{
    \eA{Then waving her hand so as to bid him quit the room, she again closed her
        eyes. Pao-yü, at the sight of her,}
}
{
}

\Verse{105}
{
    \zA{金钏儿睁眼，见是宝玉，宝玉便悄悄的笑 道：“就困的这么着？”}
}
{
    \eA{felt considerable affection for her and unable to tear himself away, so
        quietly stretching his head forward, and noticing that Madame Wang's eyes
        were shut,}
}
{
}
\ChapterImage{img/chapters/086第30-31回 椿齡畫薔癡及局外 心思踢人錯踢襲人 撕扇子作千金一笑.jpg}


\Verse{106}
{
    \zA{金钏抿嘴儿一笑，摆手叫他出去，仍合上眼。}
}
{
    \eA{he extracted from a purse, suspended about his person, one of the
        'scented-snow-for-moistening-mouth pills,'}
}
{
}

\Verse{107}
{
    \zA{宝玉见了他， 就有些恋恋不舍的，悄悄的探头瞧瞧王夫人合着眼，便自己向身边荷包里带的香雪 润津丹掏了一丸出来，向金钏儿嘴里一送，金钏儿也不睁眼，只管噙了。}
}
{
    \eA{with which it was full, and placed it on Chin Ch'uan-erh's lips. Chin
        Ch'uan-erh, however, did not open her eyes, but simply held (the pill) in
        her mouth. Pao-yü then approached her and took her hand in his. "I'll ask
        you of your mistress," he gently observed smiling, "and you and I will live
        together."}
}
{
}

\Verse{108}
{
    \zA{宝玉上来， 便拉着手，悄悄的笑道：“我和太太讨了你，咱们在一处吧？”}
}
{
    \eA{To this Chin Ch'uan-erh said not a word. "If that won't do," Pao-yü
        continued, "I'll wait for your mistress to wake and appeal to her at once."
        Chin Ch'uan-erh distended her eyes wide,}
}
{
}

\Verse{109}
{
    \zA{金钏儿不答。}
}
{
    \eA{and pushed Pao-yü off. "What's the hurry?"}
}
{
}

\Verse{110}
{
    \zA{宝玉 又道：“等太太醒了，我就说。”}
}
{
    \eA{she laughed. "'A gold hair-pin may fall into the well; but if it's yours it
        will remain yours only.'}
}
{
}

\Verse{111}
{
    \zA{金钏儿睁开眼，将宝玉一推，笑道：“你忙什么？}
}
{
    \eA{Is it possible that you don't even see the spirit of this proverb? But I'll
        tell you a smart thing. Just you go into the small court,}
}
{
}

\Verse{112}
{
    \zA{‘金簪儿掉在井里头――有你的只是有你的。’}
}
{
    \eA{on the east side, and you'll find for yourself what Mr. Chia Huau and Ts'ai
        Yun are up to!" "Let them be up to whatever they like,"}
}
{
}

\Verse{113}
{
    \zA{连这句俗语难道也不明白？}
}
{
    \eA{smiled Pao-yü, "I shall simply stick to your side!" But he then saw Madame
        Wang twist herself round,}
}
{
}

\Verse{114}
{
    \zA{我告诉 你个巧方儿：你往东小院儿里头拿环哥儿和彩云去。”}
}
{
    \eA{get up, and give a slap to Chin Ch'uan-erh on her mouth. "You mean wench!"
        she exclaimed, abusing her, while she pointed her finger at her, "it's you,
        and the like of you,}
}
{
}

\Verse{115}
{
    \zA{宝玉笑道：“谁管他的事呢! 咱们只说咱们的。”}
}
{
    \eA{who corrupt these fine young fellows with all the nice things you teach
        them!" The moment Pao-yü perceived Madame Wang rise,}
}
{
}

\Verse{116}
{
    \zA{只见王夫人翻身起来，照金钏儿脸上就打了个嘴巴，指着骂道：“下作小娼妇 儿!好好儿的爷们，都叫你们教坏了！”}
}
{
    \eA{he bolted like a streak of smoke. Chin Ch'uan-erh, meanwhile, felt half of
        her face as hot as fire, yet she did not dare utter one word of complaint.
        The various waiting-maids soon came to hear that Madame Wang had awoke and
        they rushed in in a body.}
}
{
}

\Verse{117}
{
    \zA{宝玉见王夫人起来，早一溜烟跑了。}
}
{
    \eA{"Go and tell your mother," Madame Wang thereupon said to Yü Ch'uan-erh, "to
        fetch your elder sister away."}
}
{
}

\Verse{118}
{
    \zA{这里 金钏儿半边脸火热，一声不敢言语。}
}
{
    \eA{Chin Ch'uan-erh, at these words, speedily fell on her knees. With tears in
        her eyes: "I won't venture to do it again,"}
}
{
}

\Verse{119}
{
    \zA{登时众丫头听见王夫人醒了，都忙进来。}
}
{
    \eA{she pleaded. "If you, Madame, wish to flog me, or to scold me do so at once,
        and as much as you like but don't send me away.}
}
{
}

\Verse{120}
{
    \zA{王夫 人便叫：“玉钏儿把你妈叫来!带出你姐姐去。”}
}
{
    \eA{You will thus accomplish an act of heavenly grace! I've been in attendance
        on your ladyship for about ten years, and if you now drive me away,}
}
{
}

\Verse{121}
{
    \zA{金钏儿听见，忙跪下哭道：“我 再不敢了!太太要打要骂，只管发落，别叫我出去，就是天恩了。}
}
{
    \eA{will I be able to look at any one in the face?" Though Madame Wang was a
        generous, tender-hearted person, and had at no time raised her hand to give
        a single blow to any servant-girl,}
}
{
}

\Verse{122}
{
    \zA{我跟了太太十来 年，这会子撵出去，我还见人不见人呢！”}
}
{
    \eA{she, however, when she accidentally discovered Chin Ch'uan-erh behave on
        this occasion in this barefaced manner, a manner which had all her lifetime
        been most reprehensible to her,}
}
{
}

\Verse{123}
{
    \zA{王夫人固然是个宽仁慈厚的人，从来不 曾打过丫头们一下子，今忽见金钏儿行此无耻之事，这是平生最恨的，所以气忿不 过，打了一下子，骂了几句。}
}
{
    \eA{was so overcome by passion that she gave Chin Ch'uan-erh just one slap and
        spoke to her a few sharp words. And albeit Chin Ch'uan-erh indulged in
        solicitous entreaties, she would not on any account keep her in her service.
        At length, Chin Ch'uan-erh's mother,}
}
{
}

\Verse{124}
{
    \zA{虽金钏儿苦求也不肯收留，到底叫了金钏儿的母亲白 老媳妇儿领出去了。}
}
{
    \eA{Dame Pao, was sent for to take her away. Chin Ch'uan-erh therefore had to
        conceal her disgrace, suppress her resentment, and quit the mansion. But
        without any further reference to her,}
}
{
}

\Verse{125}
{
    \zA{那金钏儿含羞忍辱的出去，不在话下。}
}
{
    \eA{we will now take up our story with Pao-yü. As soon as he saw Madame Wang
        awake, his spirits were crushed.}
}
{
}

\Verse{126}
{
    \zA{且说宝玉见王夫人醒了，自己没趣，忙进大观园来。}
}
{
    \eA{All alone he hastily made his way into the Ta Kuan garden. Here his
        attention was attracted by the ruddy sun, shining in the zenith, the shade
        of the trees extending far and wide, the song of the cicadas,}
}
{
}

\Verse{127}
{
    \zA{只见赤日当天，树阴匝地， 满耳蝉声，静无人语。}
}
{
    \eA{filling the ear; and by a perfect stillness, not even broken by the echo of
        a human voice. But the instant he got near the trellis,}
}
{
}

\Verse{128}
{
    \zA{刚到了蔷薇架，只听见有人哽噎之声。}
}
{
    \eA{with the cinnamon roses, the sound of sobs fell on his ear. Doubts and
        surmises crept into Pao-yü's mind,}
}
{
}

\Verse{129}
{
    \zA{宝玉心中疑惑，便站 住细听，果然那边架下有人。}
}
{
    \eA{so halting at once, he listened with intentness. Then actually he discerned
        some one on the off-side of the trellis. This was the fifth moon,}
}
{
}

\Verse{130}
{
    \zA{此时正是五月，那蔷薇花叶茂盛之际，宝玉悄悄的隔 着药栏一看，只见一个女孩子蹲在花下，手里拿着根别头的簪子在地下抠土，一面 悄悄的流泪。}
}
{
    \eA{the season when the flowers and foliage of the cinnamon roses were in full
        bloom. Furtively peeping through an aperture in the fence, Pao-yü saw a
        young girl squatting under the flowers and digging the ground with a
        hair-pin she held in her hand. As she dug, she silently gave way to tears.
        "Can it be possible," mused Pao-yü, "that this girl too is stupid?}
}
{
}

\Verse{131}
{
    \zA{宝玉心中想道：“难道这也是个痴丫头，又像颦儿来葬花不成？”}
}
{
    \eA{Can she also be following P'in Erh's example and come to inter flowers? Why
        if she's likewise really burying flowers," he afterwards went on to
        smilingly reflect, "this can aptly be termed:}
}
{
}

\Verse{132}
{
    \zA{因 又自笑道：“若真也葬花，可谓‘东施效颦’了，不但不为新奇，而且更是可厌。”}
}
{
    \eA{'Tung Shih tries to imitate a frown.' But not only is what she does not
        original, but it is despicable to boot. You needn't," he meant to shout out
        to the girl, at the conclusion of this train of thought,}
}
{
}

\Verse{133}
{
    \zA{想毕，便要叫那女子说：“你不用跟着林姑娘学了。”}
}
{
    \eA{"try and copy Miss Lin's example." But before the words had issued from his
        mouth, he luckily scrutinised her a second time,}
}
{
}

\Verse{134}
{
    \zA{话未出口，幸而再看时，这 女孩子面生，不是个侍儿，倒像是那十二个学戏的女孩子里头的一个，却辨不出他 是生、旦、净、丑那一个脚色来。}
}
{
    \eA{and found that the girl's features were quite unfamiliar to him, that she
        was no menial, and that she looked like one of the twelve singing maids, who
        were getting up the plays. He could not, however, make out what \_rôles\_
        she filled: scholars, girls, old men, women, or buffoons. Pao-yü quickly put
        out his tongue and stopped his mouth with his hand.}
}
{
}

\Verse{135}
{
    \zA{宝玉把舌头一伸，将口掩住，自己想道：“幸而 不曾造次。}
}
{
    \eA{"How fortunate," he inwardly soliloquised, "that I didn't make any reckless
        remark! It was all because of my inconsiderate talk on the last two
        occasions, that P'in Erh got angry with me,}
}
{
}

\Verse{136}
{
    \zA{上两回皆因造次了，颦儿也生气，宝儿也多心。}
}
{
    \eA{and that Pao-ch'ai felt hurt. And had I now given them offence also, I would
        have been in a still more awkward fix!" While wrapt in these thoughts,}
}
{
}

\Verse{137}
{
    \zA{如今再得罪了他们，越 发没意思了。”}
}
{
    \eA{he felt much annoyance at not being able to recognise who she was. But on
        further minute inspection,}
}
{
}

\Verse{138}
{
    \zA{一面想，一面又恨不认得这个是谁。}
}
{
    \eA{he noticed that this maiden, with contracted eyebrows, as beautiful as the
        hills in spring,}
}
{
}

\Verse{139}
{
    \zA{再留神细看，见这女孩子眉蹙 春山，眼颦秋水，面薄腰纤，袅袅婷婷，大有黛玉之态。}
}
{
    \eA{frowning eyes as clear as the streams in autumn, a face, with transparent
        skin, and a slim waist, was elegant and beautiful and almost the very image
        of Lin Tai-yü. Pao-yü could not, from the very first, make up his mind to
        wrench himself away.}
}
{
}

\Verse{140}
{
    \zA{宝玉早又不忍弃他而去， 只管痴看。}
}
{
    \eA{But as he stood gazing at her in a doltish mood, he realised that, although
        she was tracing on the ground with the gold hair-pin,}
}
{
}

\Verse{141}
{
    \zA{只见他虽然用金簪画地，并不是掘土埋花，竟是向土上画字。}
}
{
    \eA{she was not digging a hole to bury flowers in, but was merely delineating
        characters on the surface of the soil. Pao-yü's eyes followed the hair-pin
        from first to last, as it went up and as it came down.}
}
{
}

\Verse{142}
{
    \zA{宝玉拿眼随着簪 子的起落，一直到底，一画、一点、一勾的看了去，数一数，十八笔。}
}
{
    \eA{He watched each dash, each dot and each hook. He counted the strokes. They
        numbered eighteen. He himself then set to work and sketched with his finger
        on the palm of his hand, the lines, in their various directions,}
}
{
}

\Verse{143}
{
    \zA{自己又在手 心里拿指头按着他方才下笔的规矩写了，猜是个什么字。}
}
{
    \eA{and in the order they had been traced a few minutes back, so as to endeavour
        to guess what the character was. On completing the sketch, he discovered,
        the moment he came to reflect,}
}
{
}

\Verse{144}
{
    \zA{写成一想，原来就是个蔷 薇花的“蔷”字。}
}
{
    \eA{that it was the character "Ch'iang," in the combination, 'Ch'iang Wei,'
        representing cinnamon roses.}
}
{
}

\Verse{145}
{
    \zA{宝玉想道：“必定是他也要做诗填词，这会子见了这花，因有所 感。}
}
{
    \eA{"She too," pondered Pao-yü, "must have been bent upon writing verses, or
        supplying some line or other, and at the sight now of the flowers, the idea
        must have suggested itself to her mind.}
}
{
}

\Verse{146}
{
    \zA{或者偶成了两句，一时兴至，怕忘了，在地下画着推敲，也未可知。}
}
{
    \eA{Or it may very likely be that having spontaneously devised a couplet, she
        got suddenly elated and began, for fear it should slip from her memory, to
        trace it on the ground so as to tone the rhythm.}
}
{
}

\Verse{147}
{
    \zA{且看他底 下再写什么。”}
}
{
    \eA{Yet there's no saying. Let me see, however, what she's going to write next."}
}
{
}

\Verse{148}
{
    \zA{一面想，一面又看，只见那女孩子还在那里画呢。}
}
{
    \eA{While cogitating, he looked once more. Lo, the girl was still tracing. But
        tracing up or tracing down, it was ever the character "Ch'iang."}
}
{
}

\Verse{149}
{
    \zA{画来画去，还是 个“蔷”字；再看，还是个“蔷”字。}
}
{
    \eA{When he gazed again, it was still the self-same Ch'iang. The one inside the
        fence fell, in fact, from an early stage,}
}
{
}

\Verse{150}
{
    \zA{里面的原是早已痴了，画完一个“蔷”又画 一个“蔷”，已经画了有几十个。}
}
{
    \eA{into a foolish mood, and no sooner was one 'Ch'iang,' finished than she
        started with another; so that she had already written several tens of them.
        The one outside gazed and gazed,}
}
{
}

\Verse{151}
{
    \zA{外面的不觉也看痴了，两个眼睛珠儿只管随着簪 子动，心里却想：“这女孩子一定有什么说不出的心事，才这么个样儿。}
}
{
    \eA{until he unwittingly also got into the same foolish mood. Intent with his
        eyes upon following the movements of the pin, in his mind, he communed thus
        with his own thoughts: "This girl must, for a certainty, have something to
        say, or some unspeakable momentous secret that she goes on like this.}
}
{
}

\Verse{152}
{
    \zA{外面他既 是这个样儿，心里还不知怎么熬煎呢？}
}
{
    \eA{But if outwardly she behaves in this wise, who knows what anguish she mayn't
        suffer at heart? And yet, with a frame to all appearances so very delicate,}
}
{
}

\Verse{153}
{
    \zA{看他的模样儿这么单薄，心里那里还搁的住 熬煎呢？}
}
{
    \eA{how could she ever resist much inward anxiety! Woe is me that I'm unable to
        transfer some part of her burden on to my own shoulders!"}
}
{
}

\Verse{154}
{
    \zA{可恨我不能替你分些过来。”}
}
{
    \eA{In midsummer, cloudy and bright weather are uncertain. A few specks of
        clouds suffice to bring about rain.}
}
{
}

\Verse{155}
{
    \zA{却说伏中阴晴不定，片云可以致雨，忽然凉风过处，飒飒的落下一阵雨来。}
}
{
    \eA{Of a sudden, a cold blast swept by, and tossed about by the wind fell a
        shower of rain. Pao-yü perceived that the water trickling down the girl's
        head saturated her gauze attire in no time.}
}
{
}

\Verse{156}
{
    \zA{宝 玉看那女孩子头上往下滴水，把衣裳登时湿了。}
}
{
    \eA{"It's pouring," Pao-yü debated within himself, "and how can a frame like
        hers resist the brunt of such a squall." Unable therefore to restrain
        himself,}
}
{
}

\Verse{157}
{
    \zA{宝玉想道：“这是下雨了，他这个 身子，如何禁得骤雨一激。”}
}
{
    \eA{he vehemently shouted: "Leave off writing! See, it's pouring; you're wet
        through!" The girl caught these words, and was frightened out of her wits.}
}
{
}

\Verse{158}
{
    \zA{因此禁不住便说道：“不用写了，你看身上都湿了。”}
}
{
    \eA{Raising her head, she at once descried some one or other standing beyond the
        flowers and calling out to her: "Leave off writing. It's pouring!"}
}
{
}

\Verse{159}
{
    \zA{那女孩子听说，倒唬了一跳，抬头一看，只见花外一个人叫他“不用写了”。}
}
{
    \eA{But as Pao-yü was, firstly, of handsome appearance, and as secondly the
        luxuriant abundance of flowers and foliage screened with their boughs,
        thick-laden with leaves,}
}
{
}

\Verse{160}
{
    \zA{一则 宝玉脸面俊秀，二则花叶繁茂，上下俱被枝叶隐住，刚露着半边脸儿：那女孩子只 当也是个丫头，再不想是宝玉，因笑道：“多谢姐姐提醒了我。}
}
{
    \eA{the upper and lower part of his person, just leaving half of his countenance
        exposed to view, the maiden simply jumped at the conclusion that he must be
        a servant girl, and never for a moment dreamt that it might be Pao-yü. "Many
        thanks, sister, for recalling me to my senses,"}
}
{
}

\Verse{161}
{
    \zA{难道姐姐在外头有 什么遮雨的？”}
}
{
    \eA{she consequently smiled. "Yet is there forsooth anything outside there to
        protect you from the rain?"}
}
{
}

\Verse{162}
{
    \zA{一句提醒了宝玉，“嗳哟”了一声，才觉得浑身冰凉。}
}
{
    \eA{This single remark proved sufficient to recall Pao-yü to himself. With an
        exclamation of "Ai-yah," he at length became conscious that his whole body
        was cold as ice.}
}
{
}

\Verse{163}
{
    \zA{低头看看自 己身上，也都湿了。}
}
{
    \eA{Then drooping his head, he realised that his own person too was drenched.}
}
{
}

\Verse{164}
{
    \zA{说：“不好！”}
}
{
    \eA{"This will never do," he cried,}
}
{
}

\Verse{165}
{
    \zA{只得一气跑回怡红院去了。}
}
{
    \eA{and with one breath he had to run back into the I Hung court. His mind,
        however, continued much exercised about the girl as she had nothing to
        shelter her from the rain.}
}
{
}

\Verse{166}
{
    \zA{心里却还记挂着那 女孩子没处避雨。}
}
{
    \eA{As the next day was the dragon-boat festival, Wen Kuan and the other singing
        girls, twelve in all, were given a holiday,}
}
{
}

\Verse{167}
{
    \zA{原来明日是端阳节，那文官等十二个女孩子都放了学，进园来各处玩耍。}
}
{
    \eA{so they came into the garden and amused themselves by roaming everywhere and
        anywhere. As luck would have it, the two girls Pao-Kuan, who filled the
        \_rôle\_ of young men and Yü Kuan, who represented young women,}
}
{
}

\Verse{168}
{
    \zA{可巧 小生宝官正旦玉官两个女孩子，正在怡红院和袭人玩笑，被雨阻住，大家堵了沟， 把水积在院内，拿些绿头鸭、花？？}
}
{
    \eA{were in the I Hung court enjoying themselves with Hsi Jen, when rain set in
        and they were prevented from going back, so in a body they stopped up the
        drain to allow the water to accumulate in the yard. Then catching those that
        could be caught, and driving those that had to be driven,}
}
{
}

\Verse{169}
{
    \zA{、彩鸳鸯，捉的捉，赶的赶，缝了翅膀，放在 院内玩耍，将院门关了。}
}
{
    \eA{they laid hold of a few of the green-headed ducks, variegated marsh-birds
        and coloured mandarin-ducks, and tying their wings they let them loose in
        the court to disport themselves.}
}
{
}

\Verse{170}
{
    \zA{袭人等都在游廊上嘻笑。}
}
{
    \eA{Closing the court Hsi Jen and her playmates stood together under the
        verandah and enjoyed the fun.}
}
{
}

\Verse{171}
{
    \zA{宝玉见关着门，便用手扣门，里 面诸人只顾笑，那里听见。}
}
{
    \eA{Pao-yü therefore found the entrance shut. He gave a rap at the door. But as
        every one inside was bent upon laughing, they naturally did not catch the
        sound;}
}
{
}

\Verse{172}
{
    \zA{叫了半日，拍得门山响，里面方听见了。}
}
{
    \eA{and it was only after he had called and called, and made a noise by thumping
        at the door,}
}
{
}

\Verse{173}
{
    \zA{料着宝玉这会 子再不回来的，袭人笑道：“谁这会子叫门？}
}
{
    \eA{that they at last heard. Imagining, however, that Pao-yü could not be coming
        back at that hour, Hsi Jen shouted laughing:}
}
{
}

\Verse{174}
{
    \zA{没人开去。”}
}
{
    \eA{"who's it now knocking at the door?}
}
{
}

\Verse{175}
{
    \zA{宝玉道：“是我。”}
}
{
    \eA{There's no one to go and open."}
}
{
}

\Verse{176}
{
    \zA{麝 月道：“是宝姑娘的声音。”}
}
{
    \eA{"It's I," rejoined Pao-yü. "It's Miss Pao-ch'ai's tone of voice,"}
}
{
}

\Verse{177}
{
    \zA{晴雯道：“胡说，宝姑娘这会子做什么来？”}
}
{
    \eA{added She Yüeh. "Nonsense!" cried Ch'ing Wen. "What would Miss Pao-ch'ai
        come over to do at such an hour?"}
}
{
}

\Verse{178}
{
    \zA{袭人道： “等我隔着门缝儿瞧瞧，可开就开，别叫他淋着回去。”}
}
{
    \eA{"Let me go," chimed in Hsi Jen, "and see through the fissure in the door,
        and if we can open, we'll open; for we mustn't let her go back, wet
        through." With these words,}
}
{
}

\Verse{179}
{
    \zA{说着，便顺着游廊到门前 往外一瞧，只见宝玉淋得雨打鸡一般。}
}
{
    \eA{she came along the passage to the doorway. On looking out, she espied Pao-yü
        dripping like a chicken drenched with rain. Seeing him in this plight,}
}
{
}

\Verse{180}
{
    \zA{袭人见了，又是着忙，又是好笑，忙开了门， 笑着弯腰拍手道：“那里知道是爷回来了!你怎么大雨里跑了来？”}
}
{
    \eA{Hsi Jen felt solicitous as well as amused. With alacrity, she flung the door
        wide open, laughing so heartily that she was doubled in two. "How could I
        ever have known," she said, clapping her hands, "that you had returned, Sir!
        Yet how is it that you've run back in this heavy rain?"}
}
{
}

\Verse{181}
{
    \zA{宝玉一肚子没好气，满心里要把开门的踢几脚。}
}
{
    \eA{Pao-yü had, however, been feeling in no happy frame of mind. He had fully
        resolved within himself to administer a few kicks to the person, who came to
        open the door,}
}
{
}

\Verse{182}
{
    \zA{方开了门，并不看真是谁，还 只当是那些小丫头们，便一脚踢在肋上。}
}
{
    \eA{so as soon as it was unbarred, he did not try to make sure who it was, but
        under the presumption that it was one of the servant-girls, he raised his
        leg and give her a kick on the side.}
}
{
}

\Verse{183}
{
    \zA{袭人“嗳哟”了一声。}
}
{
    \eA{"Ai-yah!" ejaculated Hsi Jen. Pao-yü nevertheless went on to abuse.}
}
{
}

\Verse{184}
{
    \zA{宝玉还骂道：“下 流东西们，我素日担待你们得了意，一点儿也不怕，越发拿着我取笑儿了！”}
}
{
    \eA{"You mean things!" he shouted. "It's because I've always treated you so
        considerately that you don't respect me in the least! And you now go to the
        length of making a laughing-stock of me!"}
}
{
}

\Verse{185}
{
    \zA{口里 说着，一低头见是袭人哭了，方知踢错了。}
}
{
    \eA{As he spoke, he lowered his head. Then catching sight of Hsi Jen, in tears,
        he realised that he had kicked the wrong person.}
}
{
}

\Verse{186}
{
    \zA{忙笑道：“嗳哟!是你来了!踢在那里了？”}
}
{
    \eA{"Hallo!" he said, promptly smiling, "is it you who've come? Where did I kick
        you?"}
}
{
}

\Verse{187}
{
    \zA{袭人从来不曾受过一句大话儿的，今忽见宝玉生气踢了他一下子，又当着许多人， 又是羞又是气又是疼，真一时置身无地。}
}
{
    \eA{Hsi Jen had never, previous to this, received even a harsh word from him.
        When therefore she on this occasion unexpectedly saw Pao-yü gave her a kick
        in a fit of anger and, what made it worse, in the presence of so many
        people, shame, resentment, and bodily pain overpowered her and she did not,
        in fact, for a time know where to go and hide herself.}
}
{
}

\Verse{188}
{
    \zA{待要怎么样，料着宝玉未必是安心踢他， 少不得忍着说道：“没有踢着，还不换衣裳去呢！”}
}
{
    \eA{She was then about to give rein to her displeasure, but the reflection that
        Pao-yü could not have kicked her intentionally obliged her to suppress her
        indignation. "Instead of kicking," she remarked, "don't you yet go and
        change your clothes?"}
}
{
}

\Verse{189}
{
    \zA{宝玉一面进房解衣，一面笑道： “我长了这么大，头一遭儿生气打人，不想偏偏儿就碰见你了。”}
}
{
    \eA{Pao-yü walked into the room. As he did so, he smiled. "Up to the age I've
        reached," he observed, "this is the first instance on which I've ever so
        thoroughly lost control over my temper as to strike any one; and, contrary
        to all my thoughts,}
}
{
}

\Verse{190}
{
    \zA{袭人一面忍痛换 衣裳，一面笑道：“我是个起头儿的人，也不论事大事小，是好是歹，自然也该从 我起。}
}
{
    \eA{it's you that happened to come in my way?" Hsi Jen, while patiently enduring
        the pain, effected the necessary change in his attire. "I've been here from
        the very first," she simultaneously added, smilingly, "so in all things,
        whether large or small, good or bad,}
}
{
}

\Verse{191}
{
    \zA{但只是别说打了我，明日顺了手，只管打起别人来。”}
}
{
    \eA{it has naturally fallen to my share to bear the brunt. But not to say
        another word about your assault on me, why, to-morrow you'll indulge your
        hand and start beating others!"}
}
{
}

\Verse{192}
{
    \zA{宝玉道：“我才也不 是安心。”}
}
{
    \eA{"I did not strike you intentionally just now," retorted Pao-yü. "Who ever
        said,"}
}
{
}

\Verse{193}
{
    \zA{袭人道：“谁说是安心呢!素日开门关门的都是小丫头们的事，他们是 憨皮惯了的，早已恨的人牙痒痒。}
}
{
    \eA{rejoined Hsi Jen, "that you did it intentionally! It has ever been the duty
        of that tribe of servant-girls to open and shut the doors, yet they've got
        into the way of being obstinate, and have long ago become such an
        abomination that people's teeth itch to revenge themselves on them.}
}
{
}

\Verse{194}
{
    \zA{他们也没个怕惧，要是他们，踢一下子唬唬也好。}
}
{
    \eA{They don't know, besides, what fear means. So had you first assured yourself
        that it was they and given them a kick,}
}
{
}

\Verse{195}
{
    \zA{刚才是我淘气，不叫开门的。”}
}
{
    \eA{a little intimidating would have done them good. But I'm at the bottom of
        the mischief that happened just now,}
}
{
}

\Verse{196}
{
    \zA{说着，那雨已住了，宝官玉官也早去了。}
}
{
    \eA{for not calling those, upon whom it devolves, to come and open for you."
        During the course of their conversation,}
}
{
}

\Verse{197}
{
    \zA{袭人只觉肋下疼的心里发闹，晚饭也 不曾吃。}
}
{
    \eA{the rain ceased, and Pao Kuan and Yü Kuan had been able to take their leave.
        Hsi Jen, however, experienced such intense pain in her side,}
}
{
}

\Verse{198}
{
    \zA{到晚间脱了衣服，只见肋上青了碗大的一块，自己倒唬了一跳，又不好声 张。}
}
{
    \eA{and felt such inward vexation, that at supper she could not put a morsel of
        anything in her mouth. When in the evening, the time came for her to have
        her bath, she discovered, on divesting herself of her clothes,}
}
{
}

\Verse{199}
{
    \zA{一时睡下，梦中作痛，由不得“嗳哟”之声从睡中哼出。}
}
{
    \eA{a bluish bruise on her side of the size of a saucer and she was very much
        frightened. But as she could not very well say anything about it to any one,}
}
{
}

\Verse{200}
{
    \zA{宝玉虽说不是安心， 因见袭人懒懒的，心里也不安稳。}
}
{
    \eA{she presently retired to rest. But twitches of pain made her involuntarily
        moan in her dreams and groan in her sleep.}
}
{
}

\Verse{201}
{
    \zA{半夜里听见袭人“嗳哟”，便知踢重了，自己下 床来，悄悄的秉灯来照。}
}
{
    \eA{Pao-yü did, it is true, not hurt her with any malice, but when he saw Hsi
        Jen so listless and restless, and suddenly heard her groan in the course of
        the night,}
}
{
}

\Verse{202}
{
    \zA{刚到床前，只见袭人嗽了两声，吐出一口痰来，嗳哟一声， 睁眼见了宝玉，倒唬了一跳，道：“作什么？”}
}
{
    \eA{he realised how severely he must have kicked her. So getting out of bed, he
        gently seized the lantern and came over to look at her. But as soon as he
        reached the side of her bed, he perceived Hsi Jen expectorate, with a retch,
        a whole mouthful of phlegm.}
}
{
}

\Verse{203}
{
    \zA{宝玉道：“你梦里‘嗳哟’，必是 踢重了。}
}
{
    \eA{"Oh me!" she gasped, as she opened her eyes. The presence of Pao-yü startled
        her out of her wits. "What are you up to?" she asked.}
}
{
}

\Verse{204}
{
    \zA{我瞧瞧。”}
}
{
    \eA{"You groaned in your dreams,"}
}
{
}

\Verse{205}
{
    \zA{袭人道：“我头上发晕，嗓子里又腥又甜，你倒照一照地下罢。”}
}
{
    \eA{answered Pao-yü, "so I must have kicked you hard. Do let me see!" "My head
        feels giddy," said Hsi Jen. "My throat foul and sweet; throw the light on
        the floor!"}
}
{
}

\Verse{206}
{
    \zA{宝玉听说，果然持灯向地下一照，只见一口鲜血在地。}
}
{
    \eA{At these words, Pao-yü actually raised the lantern. The moment he cast the
        light below, he discerned a quantity of fresh blood on the floor.}
}
{
}

\Verse{207}
{
    \zA{宝玉慌了，只说：“了不得 了！”}
}
{
    \eA{Pao-yü was seized with consternation. "Dreadful!" was all he could say.}
}
{
}

\Verse{208}
{
    \zA{袭人见了，也就心冷了半截。}
}
{
    \eA{At the sight of the blood, Hsi Jen's heart too partly waxed cold.}
}
{
}

\Verse{209}
{
    \zA{要知端的，下回分解。}
}
{
    \eA{But, reader, the next chapter will reveal the sequel,}
}
{
}

\Verse{210}
{
    \zA{(本章完)}
}
{
    \eA{if you really have any wish to know more about them.}
}
{
}
